\documentclass[10pt,a4paper]{amsart}
\usepackage{amsmath,amssymb,amsthm,graphicx}
\usepackage{verbatim}
\newenvironment{reference}{\comment}{\endcomment}
%\newenvironment{reference}{}{}
\newenvironment{slogan}{\comment}{\endcomment}
\newenvironment{history}{\comment}{\endcomment}
% For commutative diagrams we use Xy-pic
\usepackage[all]{xy}

% We use 2cell for 2-commutative diagrams.
\xyoption{2cell}
% Theorem environments.
%
\theoremstyle{plain}
\newtheorem{theorem}[subsection]{Theorem}
\newtheorem{proposition}[subsection]{Proposition}
\newtheorem{lemma}[subsection]{Lemma}

\theoremstyle{definition}
\newtheorem{definition}[subsection]{Definition}
\newtheorem{example}[subsection]{Example}
\newtheorem{exercise}[subsection]{Exercise}
\newtheorem{situation}[subsection]{Situation}

\theoremstyle{remark}
\newtheorem{remark}[subsection]{Remark}
\newtheorem{remarks}[subsection]{Remarks}

\numberwithin{equation}{subsection}

% Macros
%
\def\lim{\mathop{\mathrm{lim}}\nolimits}
\def\colim{\mathop{\mathrm{colim}}\nolimits}
\def\Spec{\mathop{\mathrm{Spec}}}
\def\Hom{\mathop{\mathrm{Hom}}\nolimits}
\def\Ext{\mathop{\mathrm{Ext}}\nolimits}
\def\SheafHom{\mathop{\mathcal{H}\!\mathit{om}}\nolimits}
\def\SheafExt{\mathop{\mathcal{E}\!\mathit{xt}}\nolimits}
\def\Sch{\mathit{Sch}}
\def\Mor{\mathop{\mathrm{Mor}}\nolimits}
\def\Ob{\mathop{\mathrm{Ob}}\nolimits}
\def\Sh{\mathop{\mathit{Sh}}\nolimits}
\def\NL{\mathop{N\!L}\nolimits}
\def\CH{\mathop{\mathrm{CH}}\nolimits}
\def\proetale{{pro\text{-}\acute{e}tale}}
\def\etale{{\acute{e}tale}}
\def\QCoh{\mathit{QCoh}}
\def\Ker{\mathop{\mathrm{Ker}}}
\def\Im{\mathop{\mathrm{Im}}}
\def\Coker{\mathop{\mathrm{Coker}}}
\def\Coim{\mathop{\mathrm{Coim}}}

% Boxtimes
%
\DeclareMathSymbol{\boxtimes}{\mathbin}{AMSa}{"02}

%
% Macros for moduli stacks/spaces
%
\def\QCohstack{\mathcal{QC}\!\mathit{oh}}
\def\Cohstack{\mathcal{C}\!\mathit{oh}}
\def\Spacesstack{\mathcal{S}\!\mathit{paces}}
\def\Quotfunctor{\mathrm{Quot}}
\def\Hilbfunctor{\mathrm{Hilb}}
\def\Curvesstack{\mathcal{C}\!\mathit{urves}}
\def\Polarizedstack{\mathcal{P}\!\mathit{olarized}}
\def\Complexesstack{\mathcal{C}\!\mathit{omplexes}}
% \Pic is the operator that assigns to X its picard group, usage \Pic(X)
% \Picardstack_{X/B} denotes the Picard stack of X over B
% \Picardfunctor_{X/B} denotes the Picard functor of X over B
\def\Pic{\mathop{\mathrm{Pic}}\nolimits}
\def\Picardstack{\mathcal{P}\!\mathit{ic}}
\def\Picardfunctor{\mathrm{Pic}}
\def\Deformationcategory{\mathcal{D}\!\mathit{ef}}
\usepackage{cleveref}
\crefname{originallemma}{original lemma}{original lemmas}
\Crefname{originallemma}{Original lemma}{Original lemmas}
\makeatletter
\newlabel{lemma-euler-sequence}{{0FMH}{source}{}{sourcepointer-1}{}}
\newlabel{lemma-euler-sequence@cref}{{[originallemma][0][]0FMH}{[1][0][]source}{}{}{}}
\makeatother
\crefname{originallemma}{original lemma}{original lemmas}
\Crefname{originallemma}{Original lemma}{Original lemmas}
\makeatletter
\newlabel{lemma-twisted-hodge-cohomology-projective-space}{{0FUK}{source}{}{sourcepointer-2}{}}
\newlabel{lemma-twisted-hodge-cohomology-projective-space@cref}{{[originallemma][0][]0FUK}{[1][0][]source}{}{}{}}
\makeatother
\crefname{originalsection}{original section}{original sections}
\Crefname{originalsection}{Original section}{Original sections}
\makeatletter
\newlabel{section-first-chern-class}{{0FLE}{source}{}{sourcepointer-3}{}}
\newlabel{section-first-chern-class@cref}{{[originalsection][0][]0FLE}{[1][0][]source}{}{}{}}
\makeatother
\title{The Hodge cohomology of projective space is one-dimensional on the diagonal, generated by powers of the first Chern class}
\author{The Stacks Project Authors}
\date{stacks-project:derham.tex:local-snapshot-20261012; GNU FDL 1.2 or later}
\begin{document}
\maketitle
\noindent\textit{Editorial note.} All delivered mathematical text is the author's own, copied from the pinned Stacks Project chapter file derham.tex; no statement or proof step was rewritten, reordered or removed.\par\medskip
\noindent\textit{Editorial note.} Preamble provenance: the pinned chapter file opens with \textbackslash{}input\{preamble\} and carries no preamble text of its own, so the project's single Stacks-wide preamble staged at sources/stacks-derham/preamble.tex is excerpted verbatim as three spans, lines 8-12 (verbatim package and the project's own reference/slogan/history display environments), lines 14-18 (Xy-pic setup) and lines 154-217 (theorem declarations and mathematical macros). The wrapper class and the amsmath/amssymb/amsthm/graphicx base packages are supplied by the assembler. No macro was invented or redefined.\par\medskip
\noindent\textit{Editorial note.} Same-chapter pointers are located at their \textbackslash{}label lines inside the pinned chapter file; cross-chapter pointers are located in the project's own tag list staged in the snapshot (sources/stacks-chow/tags.txt), at the verbatim lines carrying their canonical tags.\par\medskip
\noindent\textit{Allowed standard prerequisites (editorial).}\begin{itemize}
\item The twisted Hodge cohomology of projective space. Same chapter derham.tex, label lemma-twisted-hodge-cohomology-projective-space (invoked at proof lines 1714, 1724 and 1756). It supplies the vanishing outside the diagonal and the freeness of rank one on the diagonal, together with the identification of the generator as a boundary in the long exact sequence.
\item The Euler sequence on projective space. Same chapter derham.tex, label lemma-euler-sequence (invoked at proof line 1723). Its short exact sequence \$0 \textbackslash{}to \textbackslash{}Omega \textbackslash{}to \textbackslash{}mathcal\{O\}(-1)\textasciicircum{}\{\textbackslash{}oplus n+1\} \textbackslash{}to \textbackslash{}mathcal\{O\} \textbackslash{}to 0\$ is what allows the generator of \$H\textasciicircum{}1(\textbackslash{}mathbf\{P\}\textasciicircum{}n\_A, \textbackslash{}Omega)\$ to be represented by an explicit cocycle.
\item The first Chern class of an invertible sheaf in the Hodge setting. Same chapter derham.tex, Section \textbackslash{}ref\{section-first-chern-class\} (invoked at proof lines 1748 and 1764). It identifies the cocycle of \$c\_1\textasciicircum{}\{Hodge\}(\textbackslash{}mathcal\{O\}(1))\$ with the logarithmic derivative cocycle, and gives the representative of its powers used in the induction.
\end{itemize}\medskip
\ifdefined\hypertarget\hypertarget{sourcepointer-1}{}\fi
\noindent\textit{Source reference.} Original lemma 0FMH: sources/stacks-derham/pinned-original-source.tex, source lines 1594--1594. Navigation pointer for \textbackslash{}ref\{lemma-euler-sequence\} used inside the counted span (line 1723); the argument that consumes it is retained in the span.\par\medskip
\ifdefined\hypertarget\hypertarget{sourcepointer-2}{}\fi
\noindent\textit{Source reference.} Original lemma 0FUK: sources/stacks-derham/pinned-original-source.tex, source lines 1644--1644. Navigation pointer for \textbackslash{}ref\{lemma-twisted-hodge-cohomology-projective-space\} used inside the counted span (line 1714); the argument that consumes it is retained in the span.\par\medskip
\ifdefined\hypertarget\hypertarget{sourcepointer-3}{}\fi
\noindent\textit{Source reference.} Original section 0FLE: sources/stacks-derham/pinned-original-source.tex, source lines 1035--1035. Navigation pointer for \textbackslash{}ref\{section-first-chern-class\} used inside the counted span (line 1748); the argument that consumes it is retained in the span.\par\medskip
\begin{lemma}
\label{lemma-hodge-cohomology-projective-space}
We have $H^q(\mathbf{P}^n_A, \Omega^p) = 0$
unless $0 \leq p = q \leq n$. For $0 \leq p \leq n$ the $A$-module
$H^p(\mathbf{P}^n_A, \Omega^p)$ free of rank $1$ with basis element
$c_1^{Hodge}(\mathcal{O}(1))^p$.
\end{lemma}

\begin{proof}
We have the vanishing and freeness by
Lemma \ref{lemma-twisted-hodge-cohomology-projective-space}.
For $p = 0$ it is certainly true that
$1 \in H^0(\mathbf{P}^n_A, \mathcal{O})$ is a generator.

\medskip\noindent
Proof for $p = 1$. Consider the short exact sequence
$$
0 \to \Omega \to \mathcal{O}(-1)^{\oplus n + 1} \to \mathcal{O} \to 0
$$
of Lemma \ref{lemma-euler-sequence}. In the proof of
Lemma \ref{lemma-twisted-hodge-cohomology-projective-space}
we have seen that the generator of $H^1(\mathbf{P}^n_A, \Omega)$
is the boundary $\xi$ of $1 \in H^0(\mathbf{P}^n_A, \mathcal{O})$.
As in the proof of Lemma \ref{lemma-euler-sequence} we will identify
$\mathcal{O}(-1)^{\oplus n + 1}$ with
$\bigoplus_{j = 0, \ldots, n} \mathcal{O}(-1)\text{d}T_j$.
Consider the open covering
$$
\mathcal{U} : 
\mathbf{P}^n_A =
\bigcup\nolimits_{i = 0, \ldots, n} D_{+}(T_i)
$$
We can lift the restriction of the global section $1$ of $\mathcal{O}$
to $U_i = D_+(T_i)$ by the section $T_i^{-1} \text{d}T_i$ of
$\bigoplus \mathcal{O}(-1)\text{d}T_j$ over $U_i$. Thus the cocyle
representing $\xi$ is given by
$$
T_{i_1}^{-1} \text{d}T_{i_1} - T_{i_0}^{-1} \text{d}T_{i_0} =
\text{d}\log(T_{i_1}/T_{i_0}) \in \Omega(U_{i_0i_1})
$$
On the other hand, for each $i$ the section $T_i$ is a trivializing
section of $\mathcal{O}(1)$ over $U_i$. Hence we see that
$f_{i_0i_1} = T_{i_1}/T_{i_0} \in \mathcal{O}^*(U_{i_0i_1})$
is the cocycle representing $\mathcal{O}(1)$ in $\Pic(\mathbf{P}^n_A)$,
see Section \ref{section-first-chern-class}.
Hence $c_1^{Hodge}(\mathcal{O}(1))$
is given by the cocycle $\text{d}\log(T_{i_1}/T_{i_0})$
which agrees with what we got for $\xi$ above.

\medskip\noindent
Proof for general $p$ by induction. The base cases $p = 0, 1$ were handled
above. Assume $p > 1$. In the proof of
Lemma \ref{lemma-twisted-hodge-cohomology-projective-space}
we have seen that the generator of $H^p(\mathbf{P}^n_A, \Omega^p)$
is the boundary of $c_1^{Hodge}(\mathcal{O}(1))^{p - 1}$
in the long exact cohomology sequence associated to
$$
0 \to \Omega^p \to \wedge^p\left(\mathcal{O}(-1)^{\oplus n + 1}\right)
\to \Omega^{p - 1} \to 0
$$
By the calculation in Section \ref{section-first-chern-class}
the cohomology class $c_1^{Hodge}(\mathcal{O}(1))^{p - 1}$
is, up to a sign, represented by the cocycle with terms
$$
\beta_{i_0 \ldots i_{p - 1}} =
\text{d}\log(T_{i_1}/T_{i_0}) \wedge
\text{d}\log(T_{i_2}/T_{i_1}) \wedge \ldots \wedge
\text{d}\log(T_{i_{p - 1}}/T_{i_{p - 2}})
$$
in $\Omega^{p - 1}(U_{i_0 \ldots i_{p - 1}})$. These
$\beta_{i_0 \ldots i_{p - 1}}$ can be lifted to the sections
$\tilde \beta_{i_0 \ldots i_{p -1}} =
T_{i_0}^{-1}\text{d}T_{i_0} \wedge \beta_{i_0 \ldots i_{p - 1}}$
of $\wedge^p(\bigoplus \mathcal{O}(-1) \text{d}T_j)$ over
$U_{i_0 \ldots i_{p - 1}}$. We conclude that the generator of
$H^p(\mathbf{P}^n_A, \Omega^p)$ is given by the cocycle whose
components are
\begin{align*}
\sum\nolimits_{a = 0}^p (-1)^a
\tilde \beta_{i_0 \ldots \hat{i_a} \ldots i_p}
& =
T_{i_1}^{-1}\text{d}T_{i_1} \wedge \beta_{i_1 \ldots i_p}
+ \sum\nolimits_{a = 1}^p (-1)^a
T_{i_0}^{-1}\text{d}T_{i_0} \wedge
\beta_{i_0 \ldots \hat{i_a} \ldots i_p} \\
& =
(T_{i_1}^{-1}\text{d}T_{i_1} - T_{i_0}^{-1}\text{d}T_{i_0}) \wedge
\beta_{i_1 \ldots i_p} +
T_{i_0}^{-1}\text{d}T_{i_0} \wedge \text{d}(\beta)_{i_0 \ldots i_p} \\
& =
\text{d}\log(T_{i_1}/T_{i_0}) \wedge \beta_{i_1 \ldots i_p}
\end{align*}
viewed as a section of $\Omega^p$ over $U_{i_0 \ldots i_p}$.
This is up to sign the same as the cocycle representing
$c_1^{Hodge}(\mathcal{O}(1))^p$ and the proof is complete.
\end{proof}

\end{document}
